\documentclass[12pt]{article}
\usepackage[utf8]{inputenc}
\usepackage[spanish]{babel}
\usepackage{geometry}
\usepackage{hyperref}
\usepackage{booktabs}
\usepackage{xcolor}
\usepackage{longtable}
\usepackage{array}

\geometry{a4paper, margin=2.5cm}

\definecolor{accentblue}{HTML}{1A3A5C}
\definecolor{accentteal}{HTML}{0D7377}
\definecolor{midblue}{HTML}{2E75B6}

\title{Memoria Técnica: Algorand (ALGO): Análisis Experto de la Red Blockchain}
\author{Imad Habib Jali (901421) \and Jorge Bellido Lobera (903080)}
\date{}

\begin{document}

\maketitle

\section*{Introducción}

Algorand es una red blockchain de tercera generación diseñada desde sus cimientos para superar las limitaciones que históricamente han afectado a las blockchains de primera y segunda generación. Concebida por Silvio Micali (Premio Turing 2012), la red busca resolver el "trilema blockchain" de forma matemática y robusta.

\section{Qué es Algorand}

\subsection{Origen y visión fundacional}
El proyecto fue concebido en 2017 por Silvio Micali, catedrático del MIT. Micali partió de la premisa de que el trilema blockchain (seguridad, escalabilidad y descentralización) podía resolverse mediante un mecanismo de consenso radicalmente diferente, basado en el rigor académico.

\subsection{Definición y propósito}
Algorand es una blockchain pública, permissionless y descentralizada, orientada a aplicaciones de alto rendimiento en finanzas descentralizadas (DeFi), CBDCs y activos digitales. Su token nativo es ALGO.

\subsection{Ecosistema institucional}
La Fundación Algorand (Singapur) gestiona la gobernanza y el ecosistema, mientras que Algorand Inc. (Boston) lidera el desarrollo técnico y la adopción empresarial.

\section{Cómo funciona la red blockchain de Algorand}

\subsection{Pure Proof of Stake (PPoS)}
La innovación central es el Pure Proof of Stake (PPoS). El proceso se desarrolla en tres fases:
\begin{enumerate}
    \item \textbf{Selección del líder:} Mediante Función Aleatoria Verificable (VRF), se selecciona aleatoriamente un líder de bloque.
    \item \textbf{Votación del comité:} Un comité seleccionado por VRF vota la aprobación del bloque.
    \item \textbf{Certificación y difusión:} El bloque se añade a la cadena con finalidad definitiva en menos de 5 segundos.
\end{enumerate}

\subsection{Estructura en dos capas}
\begin{itemize}
    \item \textbf{Capa 1 (Layer 1):} Aloja transacciones básicas, ASAs y contratos inteligentes estándar. Optimización para velocidad.
    \item \textbf{Capa 2 (Layer 2):} Maneja lógica compleja fuera de la cadena, reconciliando estados periódicamente con la Capa 1.
\end{itemize}

\subsection{Smart Contracts (ASC1) y el AVM}
Los contratos se ejecutan en la Algorand Virtual Machine (AVM). Se dividen en \textit{stateful} (con estado) y \textit{smart signatures} (sin estado). Se programan principalmente en PyTeal o TEAL.

\subsection{Tokenomics y sostenibilidad}
Suministro fijo de 10.000 millones de ALGO. Comisiones extremadamente bajas (~0.001 ALGO). Algorand es \textit{carbon negative} mediante acuerdos de compensación de carbono.

\section{Diferencias con Bitcoin y Ethereum}

\begingroup
\centering
\begin{longtable}{|p{4cm}|p{3.5cm}|p{3.5cm}|p{3.5cm}|}
\caption{Tabla comparativa integral} \\
\hline
\textbf{Criterio} & \textbf{Bitcoin} & \textbf{Ethereum} & \textbf{Algorand} \\ \hline
\endhead
Año de creación & 2009 & 2015 & 2017 (mainnet 2019) \\ \hline
Mecanismo de consenso & Proof of Work (PoW) & Proof of Stake (PoS) desde 2022 & Pure Proof of Stake (PPoS) \\ \hline
Velocidad (TPS) & $\sim$7 tx/seg & $\sim$15-30 tx/seg & $\sim$1.000-2.000 tx/seg \\ \hline
Finalidad de tx & $\sim$60 min (6 bloques) & $\sim$12-15 min & $<$ 5 segundos \\ \hline
Smart Contracts & No (limitado) & Sí (EVM/Solidity) & Sí (TEAL/PyTeal/AVM) \\ \hline
Consumo energético & Muy alto & Bajo (PoS) & Muy bajo (PPoS) \\ \hline
Huella de carbono & Negativa (alta) & Neutral/baja & Carbon negative \\ \hline
Comisiones (fees) & Variables, altas & Variables, altas & Muy bajas ($\sim$0.001 ALGO) \\ \hline
Suministro máximo & 21 millones BTC & Sin límite definido & 10.000 millones ALGO \\ \hline
Staking mínimo & N/A & 32 ETH & 1 ALGO \\ \hline
Gobernanza on-chain & No directamente & Parcial (EIPs) & Sí (90 días) \\ \hline
CBDCs adoptadas & No & Exploraciones & Marshall Islands, +16 proy. \\ \hline
Bifurcaciones (forks) & Posibles & Posibles & Imposibles por diseño \\ \hline
\end{longtable}
\endgroup

\section{Escenario de Uso: Tokenización de Activos Agrícolas}

\subsection{Contexto y justificación}
Se propone la tokenización de una cooperativa olivarera ("AgroSur S.L.") para democratizar la inversión y automatizar la gestión mediante contratos inteligentes.

\subsection{Desarrollo técnico paso a paso}
\begin{enumerate}
    \item \textbf{Tokenización:} Creación de tokens AGSUR mediante protocolo ASA (1 token = 1 $m^2$).
    \item \textbf{Smart Contract de gobernanza:} Reglas de dividendos y votación programadas en PyTeal.
    \item \textbf{Captación y Opt-In:} Inversores adquieren tokens y realizan el Opt-In al contrato.
    \item \textbf{Operación de campaña:} Gastos auditables y votaciones on-chain.
    \item \textbf{Distribución automática:} Pago de dividendos mediante Atomic Transfers al finalizar la campaña.
    \item \textbf{Auditoría total:} Registro inmutable y transparente en el explorador de bloques.
\end{enumerate}

\section*{Conclusiones}
Algorand representa una evolución técnica fundamentada en ciencia criptográfica. Su mecanismo PPoS ofrece una combinación única de velocidad, seguridad y descentralización. Frente a Bitcoin y Ethereum, aporta finalidad inmediata, costes despreciables y una arquitectura institucionalmente compatible y sostenible.

\end{document}
